\subsection{Absolute Value Review} 
\begin{definition}[Absolute Value]
The \term{absolute value} of a real number $x$, written $\Abs{x}$, is the distance from $x$ to $0$.
\end{definition}
To take the absolute value of a number, we ignore the sign and make it \makeblank[0.75in].
\begin{example}
Take each absolute value.
\begin{itemize}
\item $\Abs{-2}=\makeblanksmall$
\item $\Abs{5}=\makeblanksmall$
\item $\Abs{-4.7}=\makeblanksmall$
\item $\Abs{\frac{7}{3}}=\makeblanksmall$
\end{itemize}
\end{example}
\noindent In order of operations, the absolute value acts as a grouping symbol. That is, we:
\begin{enumerate}[1.]
\item Evaluate the expression inside the absolute value.
\item Take the absolute value.
\item Evaluate the remaining operations.
\end{enumerate}
\begin{example} Evaluate.
\begin{lpic}{}{5}
\regtextleft{1}{-1}{$-\frac{2}{3}\Abs{7-2\cdot 5}+8={}$};
\end{lpic}
\end{example}
